% !TeX root = ../navier-stokes-manuscript.tex
% ============================================================
% 2.1 Critical Height at a Finite Endpoint
% ============================================================

\subsection{Critical Height at a Finite Endpoint}\label{sub:CriticalHeightAtAFiniteEndpoint}

Suppose, for the sake of following the singularity threat, that a classical solution is smooth on \([0,T_*)\) and that the alleged maximal time \(T_*\) is finite.
The Blown picture from the introduction allows the first unbounded quantity to be a derivative, a pressure response, or another finite-depth readout of the fluid.
The velocity itself does not have to be the first quantity that escapes.
The first analytic question asks which measurement continues to see the fluid when its activity moves toward a smaller spatial scale.

The \NavierStokes\ equation supplies that measurement through its scaling.
The exact dilation argument belongs to the whole-space branch \(\Omega=\mathbb R^3\), which is also the branch carried by the formal closure below; the periodic problem has no global spatial dilation symmetry.
If \((u,p)\) solves the equation with viscosity \(\nu>0\), then for every \(\lambda>0\),
\[
  u_\lambda(x,t)
  :=\lambda u(\lambda x,\lambda^2t),
  \qquad
  p_\lambda(x,t)
  :=\lambda^2p(\lambda x,\lambda^2t)
\]
solves the same equation with the same viscosity on the rescaled time interval.
The factor \(\lambda\) changes the velocity amplitude, the argument \(\lambda x\) compresses length by \(1/\lambda\), and the argument \(\lambda^2t\) compresses time by \(1/\lambda^2\).
This is the parabolic relation between one temporal step and two spatial steps that will reappear inside the derivative tower.

Let \(\Lambda=(-\Delta)^{1/2}\), and measure spatial height with the homogeneous Sobolev norm
\[
  \norm{v}_{\dot H^s}^2
  :=
  \int_{\R^3}|\xi|^{2s}|\widehat v(\xi)|^2\,d\xi.
\]
The Fourier transform of the rescaled velocity is
\[
  \widehat{u_\lambda}(\xi,t)
  =
  \lambda^{-2}\widehat u(\xi/\lambda,\lambda^2t).
\]
Substitution into the Sobolev norm and the change of variables \(\eta=\xi/\lambda\) give
\begin{align*}
  \norm{u_\lambda(t)}_{\dot H^s}^{2}
  &={}
  \int_{\R^3}
  |\xi|^{2s}\lambda^{-4}
  \left|\widehat u(\xi/\lambda,\lambda^2t)\right|^2
  \,d\xi
  \\
  &={}
  \lambda^{2s-1}
  \int_{\R^3}
  |\eta|^{2s}
  \left|\widehat u(\eta,\lambda^2t)\right|^2
  \,d\eta.
\end{align*}
Thus the exact scaling law is
\[
  \norm{u_\lambda(t)}_{\dot H^s}
  =
  \lambda^{s-1/2}
  \norm{u(\lambda^2t)}_{\dot H^s}.
\]
At \(s=1/2\), the exponent vanishes.
This earns the critical height
\[
  H(t)
  :=
  \frac12\norm{u(t)}_{\dot H^{1/2}}^2
  =
  \frac12\norm{\Lambda^{1/2}u(t)}_2^2.
\]
Exact rescaling can move a configuration to finer length and shorter time while leaving this height unchanged.

The kinetic-energy scale behaves differently.
Putting \(s=0\) in the same calculation gives
\[
  \norm{u_\lambda(t)}_2^2
  =
  \lambda^{-1}
  \norm{u(\lambda^2t)}_2^2.
\]
The rescaled kinetic energy becomes smaller as the spatial concentration becomes finer.
A global \(L^2\) energy ceiling can coexist with activity that remains visible at the critical height.
This is why finite energy alone cannot settle a critical concentration scenario.

The Zeno picture makes the finite-time pressure concrete.
Take a sequence of length scales \(\ell_k=2^{-k}\).
Their parabolic response times have size \(\ell_k^2=4^{-k}\), and hence
\[
  \sum_{k=0}^{\infty}\ell_k^2
  =
  \sum_{k=0}^{\infty}4^{-k}
  <\infty.
\]
Infinitely many finer-scale episodes are dimensionally capable of accumulating before a finite time.
The scaling law also allows their critical height to remain visible while their kinetic-energy scale decreases.
This calculation describes the geometry of a possible cascade; it does not assert that the \NavierStokes\ equation realizes one.
It tells us what an exclusion argument has to follow if the threatened failure keeps moving beyond every previously inspected scale.

Scale alone still records only where the activity sits.
Smoothness also asks how the fluid responds while that activity changes.
The earlier velocity--acceleration--jerk picture now acquires a precise finite-time consequence.
